\section{The Primitive and the \CoDE Pipeline}\label{sec:approach}
\begin{definition}[Counterfactual explanation]\label{def:cfe}
Let $P$ be a function, $T$ a set of test inputs, and $m$ a single-site syntactic intervention yielding $P_m$. Write
$P(x)\simeq P_m(x)$ when both return equal values or raise the same exception type. The \emph{executed difference} is
$\Delta_T(P,m)=\{x\in T : P(x)\not\simeq P_m(x)\}$. A counterfactual explanation is $E=(m,A,d,x^\star,y,y_m)$: the
intervention, the affected-input claim $A$, a statement $d$ of how the result changes, and a witness $x^\star$ with claimed
outcomes $y,y_m$. $E$ is \emph{witness-valid} iff $P(x^\star)\not\simeq P_m(x^\star)$, $P(x^\star)\simeq y$ and
$P_m(x^\star)\simeq y_m$. If $A$ is given as a predicate $a:T\to\{0,1\}$, $E$ is \emph{affected-exact on $T$} iff
$a(x)=[x\in\Delta_T]$ for all $x\in T$.
\end{definition}
The checks cover the witness and predicate. The prose $d$ is unchecked in the execution loop. A deterministic template can copy a
verified witness from $\Delta_T$ by construction. This leaves open whether an explanation can identify new examples
or characterize the affected inputs. The outcome comparison uses float tolerances (relative $10^{-7}$, absolute $10^{-6}$)
and compares list and tuple contents, with set comparison for designated tasks. Exceptions match by type, so changing
only a \texttt{ValueError} message leaves $\Delta_T$ unchanged.

\textbf{Pipeline.} \CoDE (Counterfactual Differential Explanation) has four stages.
\emph{(1) Intervene.} AST operators propose edits in the target function: negate an \texttt{if}/\texttt{while} condition;
remove an \texttt{if} whose body is a single return/raise/break/continue; swap a comparison
($<\!\leftrightarrow\!\le$, $>\!\leftrightarrow\!\ge$, $==\!\leftrightarrow\!\neq$); swap an arithmetic operator
($+\!\leftrightarrow\!-$, $*\!\leftrightarrow\!/\!/$, $\%\!\to\!/\!/$, $/\!\to\!*$, $**\!\to\!*$); shift an integer constant
in a comparison, \texttt{range} or slice by $\pm1$; change a default argument. Candidates are tried in a seeded,
operator-balanced order.
\emph{(2) Execute.} $P$ and each candidate run on $T$ in a separate \texttt{python -I} process (empty working directory,
minimal environment, recursion clamp 1500, per-input timeouts; process isolation, see \S\ref{sec:agenda}). The first candidate
with $\Delta_T\neq\emptyset$ is kept; equivalent, non-terminating and crash-everywhere candidates are counted and skipped.
\emph{(3) Articulate.} An LLM receives $P$ and the diff and returns $(A,d,x^\star,y,y_m)$ as JSON, and, in a separate
call, $A$ as a Python predicate \texttt{affected(...)} (at most 12 lines, told not to re-implement either version).
In \emph{Ungrounded} (U), the prompt contains only $P$, $P_m$, and schema instructions. In \emph{Grounded} (G), the prompt additionally provides execution evidence: the diff ratio $|\Delta_T|/|T|$, the 3 shortest differing inputs with both executed outputs, and 2 agreeing inputs, constrained by: \emph{``Do NOT copy any shown input; identify a novel witness.''}
\emph{(4) Check.} Generated witnesses fall into four disjoint categories: (i)~\emph{valid} (causes divergence and both claimed outcomes match execution); (ii)~\emph{outcome mismatch} (causes divergence, but claimed $y$ or $y_m$ is wrong); (iii)~\emph{non-differing} ($P(x^\star)\simeq P_m(x^\star)$); and (iv)~\emph{malformed} (argument syntax error or uncaught exception in both versions). The check is lenient: both readings of an argument list for one-parameter functions are tried, and untyped ``raises'' matches any exception. Predicates pass an AST filter (no imports beyond pure math, no \texttt{open}/\texttt{eval}/\texttt{exec}) and run in the isolated runner on $T$; unparseable or non-Boolean outputs are \emph{unscorable}.

Fig.~\ref{fig:teaser} is a real case (display simplified: no docstring, entry function renamed \texttt{f}, generators
elided; ``14B'' is \GenModel, ``Haiku'' \StrongModel; ungrounded).

